\section{Related Work}
\label{sec:related}

\para{Linear representations in LLMs.}
The observation that concepts correspond to directions in neural network activation spaces dates back to word embeddings~\citep{mikolov2013efficient} and has been formalized for LLMs by \citet{park2024linear}, who define three notions of linear representation (subspace, measurement, intervention) and unify them through a \emph{causal inner product}.
They show that causally separable concepts yield orthogonal directions, and demonstrate block-diagonal structure among related concepts in LLaMA-2-7B.
\citet{park2024geometry} extend this framework to categorical and hierarchical concepts, proving that hierarchical relationships decompose into orthogonal sub-directions.
We build on their direction extraction methodology and causal inner product, but focus on a question they do not address: whether independently extracted directions compose under addition.

\para{Steering vectors and representation engineering.}
\citet{zou2023representation} introduce representation engineering, extracting ``reading vectors'' for high-level behavioral concepts (honesty, fairness) via contrastive pairs and using them as steering vectors.
\citet{panickssery2023steering} refine this approach with Contrastive Activation Addition for behavioral concepts like sycophancy.
Both lines of work demonstrate that steering with single directions is effective, but neither systematically tests multi-direction composition.
Our functional steering experiments directly address this gap, testing whether composed directions $\vd_A + \vd_B$ preserve both effects.

\para{Function vectors and composition.}
Most directly relevant to our work, \citet{todd2024function} discover that in-context learning tasks are encoded as compact function vectors extractable from causal attention heads.
They show that function vectors obey vector algebra for some task compositions (\eg $\text{Last-Capital} \approx \text{Last-Copy} + \text{First-Capital} - \text{First-Copy}$) but fail for others.
Our study differs in three ways: (1)~we test composition across a much broader taxonomy of 25~concept types rather than a handful of composed tasks, (2)~we use the unembedding space rather than intermediate attention heads, and (3)~we introduce quantitative metrics (CFS, interference scores) that measure \emph{degree} of compositionality rather than treating it as binary.

\para{Format dependence of directions.}
\citet{opielka2026causality} show that function vectors extracted via activation patching are not format-invariant: the same concept extracted from open-ended vs.\ multiple-choice formats yields nearly orthogonal directions.
They introduce \emph{concept vectors} via representational similarity analysis that are format-invariant.
This finding underscores that ``linear direction'' is not a monolithic notion --- the extraction method and context matter.
Our use of the unembedding matrix, which is a fixed model parameter rather than a context-dependent activation, sidesteps format dependence but is limited to concepts expressible through single-token contrasts.

\para{Superposition and interference.}
\citet{elhage2022superposition} demonstrate in toy models that networks pack more features than they have dimensions by encoding them as nearly-orthogonal directions in a lower-dimensional space.
This \emph{superposition} introduces interference: activating one direction partially activates others.
\citet{bereska2025superposition} provide an information-theoretic framework connecting superposition to adversarial vulnerability, and \citet{hanni2024computational} study computational superposition during inference.
Superposition is a key mechanism by which composition could fail, and our interference metric directly measures this effect.
Our finding that lexicographic relations fail to form coherent directions is consistent with these relations being encoded in highly superposed, distributed representations.
